% 原分章PDF119页，全部7道习题；120页仅出版社标识。
\Needspace{6\baselineskip}
\section*{习题}
\addcontentsline{toc}{section}{习题}
\markright{习题}

\begin{enumerate}[label=\textbf{习题\arabic*.},leftmargin=5em]
\item 考虑结构方程模型
\[
\begin{aligned}
X&=\varepsilon_x,\\
Y&=X^2+\varepsilon_y,
\end{aligned}
\]
其中 $\varepsilon_x,\varepsilon_y$ 均服从正态分布 $N(0,1)$。若将 $X$ 设为3，求效应 $Y$ 的分布。

\item 令 $f(x)=\frac1{2b}\exp\{-|x-\mu|/b\}$ 为Laplace分布的密度函数。证明其微分熵为
\[
S(l)=1+\log(2b).
\]

\item 令 $f(x)=x^{\alpha-1}(1-x)^{\beta-1}/B(\alpha,\beta)$ 为Beta分布的密度函数。证明其微分熵为
\[
\begin{aligned}
S(b)={}&\log B(\alpha,\beta)-(\alpha-1)[\psi(\alpha)-\psi(\alpha+\beta)]\\
&{}-(\beta-1)[\psi(\beta)-\psi(\alpha+\beta)],
\end{aligned}
\]
其中 $\psi$ 为双伽马函数。

\item 令 $f(x)$ 为Beta分布，$g(x)$ 为 $[0,1]$ 上的均匀分布。计算K–L距离 $D(f\parallel g)$。

\item 令
\[
g_2(x)=D(P_{Y\mid x}\parallel\vec P(y)
 =\int\log\frac{P(y\mid x)}{\vec P(y)}P(y\mid x)\,dy.
\]
讨论 $\operatorname{cov}_{U_x}\bigl(g_2,P_X/U_X\bigr)=0$ 的含义。

\item 证明：若 $X,Y$ 均为一维，则
\[
\Delta_{X\to Y}=0,\qquad\text{且}\quad\Delta_{Y\to X}=0.
\]

\item 证明
\[
\Delta_{X\to Y}+\Delta_{Y\to X}
 =-\log(1-\operatorname{cov}(Z,1/Z))+\log\frac nm.
\]
\end{enumerate}
